\section{Introduction}\label{sec:introduction}

Throughout, $\Bbbk$ is an algebraically closed field of characteristic
zero.  Formal-algebraic constructions are made over $\Bbbk$; oscillatory
integrals and Saito theory are understood after base change to $\CC$.
The enumerative coefficient rings are defined over $\QQ$ and are implicitly
base-changed to $\Bbbk$ when used in the root-stack construction.

\subsection{From boundary-ray continuation to an affine scattering problem}

The Gross--Siebert programme constructs mirrors by placing wall-crossing
automorphisms on an integral-affine base with singularities and reading off
broken lines and theta functions
\cite{grosshackingkeel2015,grossetal2018canonical,grosssiebert2022}.  What it
requires as input is geometric: a base, its singular locus, a polyhedral
decomposition, and a polarization.  This paper asks whether the wall crossing
of an open Landau--Ginzburg theory supplies that input, and settles the
question for one model.

Gross, Kelly, and Tessler define genus-zero open FJRW invariants for
\[
 (W,G)=\bigl(x^r+y^s,\mu_r\times\mu_s\bigr)
\]
as relative Euler numbers on moduli spaces of $W$-spin disks
\cite{grossetal2022open}.  Their invariants depend on canonical
multisections, and changes of those boundary conditions act through a formal
wall-crossing group closely related to the tropical vertex group.  This
action has the algebra of a scattering diagram but not, by itself, an
integral-affine realization: there is no affine manifold, singular locus, or
polyhedral decomposition on which the walls have been placed.

The companion paper \cite{wuebben2026boundary} isolates the gap precisely.

\begin{importedthm}[Boundary-ray factorization \cite{wuebben2026boundary}]
\label{thm:companion}
Fix a finite set of deformation monomials closed under taking divisors.  After
discarding every other monomial, the change between the two canonical GKT
boundary conditions at the limiting boundary rays factors uniquely into
wall-crossing operations ordered by slope.  The slope of a factor is the ray
spanned by the pair of boundary-marking counts of the corresponding critical
$W$-spin graph.  Each factor is realized by a homotopy of canonical boundary
conditions.  Radial motion in the punctured first quadrant of marking counts
acts trivially, so the product is transport from one limiting ray to the other
and not monodromy around the origin.  After the degree-zero torus action is
included, the two endpoint boundary conditions and their transition are
unique in the finite quotient.  The nonidentity factors admit successive
finite homotopies in increasing-slope order.
\end{importedthm}

This result supplies an ordered family of wall-crossing operations, but it
does not supply the integral-affine surface on which their rays should lie or
any affine monodromy.  The present paper constructs that missing geometry.

We resolve this for the parabolic quartic
\[
 W=x^4+y^4,
\]
the parabolic singularity $X_9$, whose central charge is exactly one.  The
suspension of its marginal pencil is an anticanonical elliptic curve in
$\PP(1,1,2)$, and the toric-boundary degeneration of that pair gives the
disk-shaped dual intersection complex of the Landau--Ginzburg construction of
Carl--Pumperla--Siebert \cite{carletal2024}; we call it the \emph{CPS
complex}.  It is the natural candidate base.  The first result of this paper
is that it cannot carry the descendent lattice, and that three further
standard choices fail with it: a finite
root stack of coefficients, a scalar Hamiltonian algebra, and a theta basis
indexed by one lattice.  Proposition~\ref{thm:intro-obstructions} collects the
five obstructions, each proved in the body.

\subsection{Initial Hamiltonians}

The affine and scattering constructions begin with two elements of the GKT
critical Lie algebra.  Write $M=\ZZ^2$ for the toric character lattice in the
chosen toric chart.  A
critical graph of \cite{grossetal2022open}
is supported on a nonempty set $J$ of boundary markings, each carrying a
descendent order; we call the datum a \emph{singleton} when $|J|=1$, and \emph{primary} when no marking carries a
positive descendent order.  The two first-order singleton generators of
$x^4+y^4$ are, writing $q^\perp=(q_2,-q_1)$, $\partial_v$ for the invariant
logarithmic derivation in direction $v$, and $\kappa=(1,1)$ for the shift
introduced below, in the shifted notation
$X^\kappa_q=z^{q-\kappa}\partial_{q^\perp}$ of
Section~\ref{sec:jacobi-theta},
\[
 \mathsf A=-\tfrac12t_1X^\kappa_{(2,1)},
 \qquad
 \mathsf B=-\tfrac12t_2X^\kappa_{(1,2)},
\]
the coefficient $-\tfrac12$ being the normalization borne by the GKT
invariants themselves.  The integer pairs $(2,1)$ and $(1,2)$ label the
generators and do lie in $M$.  Their wall supports are obtained by the
cofactor map $F$ in \eqref{eq:intro-F}:
\[
 F(2,1)=(1/4,-1/2),\qquad F(1,2)=(-1/4,-1).
\]
These support vectors do not lie in $M$; their integral span with $M$ is the
index-eight overlattice $\widetilde M$.  This is the lattice mismatch that
produces the annular cover of Section~\ref{sec:affine-kummer}.  We call
$\mathsf A$ and $\mathsf B$ the \emph{singleton
Hamiltonians}, and we call their deformation parameters $t_1,t_2$ the
\emph{singleton parameters}; the ideal $I=(t_1,t_2)$ is the \emph{singleton
ideal}.  These two elements are the only enumerative input to
Sections~\ref{sec:affine-kummer}--\ref{sec:compactification}.  We call
$X_{0,0}=x\partial_x-y\partial_y$ the \emph{degree-zero} generator; the word
\emph{marginal} is reserved throughout for the charge-one deformation
direction $x^2y^2$, the modulus of $X_9$.  An element of $\mu_4\times\mu_4$
is \emph{narrow} when its fixed locus in $\CC^2$ is the origin; the nine
narrow elements are $(\zeta^{a+1},\zeta^{b+1})$ with $0\leq a,b\leq2$, matching
the monomial basis $x^ay^b$ of the Jacobian ring $\Jac(x^4+y^4)$.

\subsection{Coefficient system and volume line}

Two consequences of the lattice mismatch determine the coefficient system.
Let a \emph{finite transport tree} be a finite rooted tree whose vertices are
chambers of the annular wall decomposition and whose oriented edges are
reduced paths between adjacent chambers.  For such a tree $T$ and an integer
$N$, adjoining the finitely many roots required to lift every slab and wall
map along $T$, and then reducing modulo $I^N$, gives a finite root stack
$\mathcal R_{T,N}$.  Inclusion of trees, refinement of paths, and reduction of
$N$ give transition maps.  We use the inverse system
\begin{equation}
 \{\mathcal R_{T,N}\}_{(T,N)}
 \label{eq:intro-root-system}
\end{equation}
over this cofiltered category; no geometric inverse limit is asserted.  It is
formally analogous to the finite stages of an infinite root stack
\cite{talpovistoli2018}, but the indexing category here records path transport
rather than divisibility alone.

The second uses a twisted volume form.  In a developed torus chart put
\[
 \kappa=(1,1),\qquad
 \Omega_\kappa=z^\kappa\,d\log z_1\wedge d\log z_2.
\]
$\kappa$ denotes the shift in the chosen developed chart and is transported
with that chart; it is not a global flat section of the character local
system.  A lifted slab transition preserves
$d\log z_1\wedge d\log z_2$ but rescales
$\Omega_\kappa$ by a root-valued unit, so the local Hamiltonian brackets glue
on a line bundle rather than on the structure sheaf.  The resulting object is
the Jacobi--Kirillov bracket on the \emph{volume line} $\Lcal_\kappa$;
equivalently, the complement of the zero section in $\Lcal_\kappa^\vee$ carries
a homogeneous Poisson structure, the Poissonization of the Jacobi bracket
\cite{vitaglianowade2016}.  Saito's primitive form is a different object and
enters only in Section~\ref{sec:primary-frobenius}.

\subsection{Annular monodromy and its consequences}

Two mechanisms produce the departures from the standard finite-type
situation.  The first descendent exponents force an overlattice
$\widetilde M\supset M$ of index eight; only one of the three affine
monodromies of the CPS complex preserves it; the double cover branched at the
two offending singular points is an annulus; and the holonomy $K$ of that annulus about
its core has trace $18$, hence is hyperbolic.  Repeated slab shear makes the
transported root orbit infinite and forces the inverse system
\eqref{eq:intro-root-system}.
Hyperbolicity makes every nonzero exponent orbit infinite and forces the theta
functions to form a local system rather than a basis indexed by one lattice.

\subsection{Basic notation}

We use the standard scattering-diagram terminology of
\cite{grosssiebert2022}.  In particular, slabs carry the gluing functions on
codimension-one cells, whereas walls are outgoing.  The cofactor map below
determines the supports of the two initial Hamiltonians.

Let $M\simeq\ZZ^2$ be the character lattice of the quartic toric model,
$\varphi$ its polarization, and, in the cofactor chart of
Section~\ref{sec:affine-kummer}, put
\begin{equation}\label{eq:intro-F}
 F=\frac14\begin{pmatrix}1&-1\\0&-2\end{pmatrix},
 \qquad \widetilde M=F(\ZZ^2),
 \qquad D=\widetilde M/M\simeq\ZZ/2\oplus\ZZ/4.
\end{equation}
Let $\overline\Acal$ be the completed parity cover,
$\Acal=\operatorname{Int}(\overline\Acal)$, and let
$\Acal^\circ$ be the complement of its four affine singularities.  We write
$t_{a,b}$, $0\leq a,b\leq2$, for the primary deformation parameters dual to
the monomial basis $x^ay^b$ of $\Jac(x^4+y^4)$, which index the nine narrow
sectors, and reserve $t_1,t_2$ for the singleton parameters of \S1.2.  The linear holonomy of $\Acal$ about its
core, in the basis of $\widetilde M$ fixed by $F$ and with the orientation of
Appendix~\ref{app:calculations}, is
\begin{equation}\label{eq:intro-K}
 K=\begin{pmatrix}-1&2\\-10&19\end{pmatrix},
 \qquad \operatorname{tr}K=18.
\end{equation}

\subsection{Obstructions to the disk-shaped affine base}

\begin{proposition}[Obstructions to simpler models]\label{thm:intro-obstructions}
The following hold for the open FJRW wall crossing of $(x^4+y^4,\mu_4^2)$.
\begin{enumerate}
\item\label{it:obs-lattice} \textup{(No model on the disk.)}  The first
descendent Hamiltonians do not act integrally on $M$, which has index eight
in $\widetilde M$.  Of the three affine monodromies of the disk-shaped
Carl--Pumperla--Siebert complex, exactly one preserves $\widetilde M$.  The
smallest cover on which $\widetilde M$ is preserved completes to an annulus.
Thus the enlarged lattice defines no model on that complex or on an ordinary
root-stack refinement with the same coarse disk.
\item\label{it:obs-root} \textup{(No finite root stack.)}  No finite
algebraic Kummer extension containing the initial root data is stable under
the slab shear.  If every transported root and its automorphisms are
retained, the resulting coefficient object is not a formal scheme, a finite
root stack, or an algebraic stack whose diagonal is locally of finite type.
\item\label{it:obs-scalar} \textup{(No scalar bracket.)}  The scalar
$\kappa$-twisted Hamiltonian algebra containing the two singleton fields is
not preserved under conjugation by the lifted slab transitions.  Hence its
local brackets do not glue to a bracket on the structure sheaf.
\item\label{it:obs-theta} \textup{(No global theta basis.)}  No nonzero
$m\in\widetilde M$ has finite $K$-orbit, and the formal $K$-orbit sum of a
theta section does not converge $I$-adically.  There is no theta basis
indexed by one copy of $\widetilde M$.
\item\label{it:obs-regular} \textup{(No ordinary recovery of the quartic
family, and no regular minimal presentation.)}  The annular cover is not the
dual intersection complex of a distinguished toric degeneration of del Pezzo
surfaces with simple singularities, proper generic fiber, and smooth
anticanonical divisor.  Moreover some finite root-stack presentation has singular atlas
and is not log regular.
\end{enumerate}
\end{proposition}

The five parts are proved, in order, by
Proposition~\ref{prop:integral-cofactor} with
Theorem~\ref{thm:parity-cover}; Theorems~\ref{thm:infinite-root-orbit} and
\ref{thm:forced-category}; Proposition~\ref{prop:no-scalar-frame};
Theorem~\ref{thm:theta-local-system}; and
Propositions~\ref{prop:no-CPS-degeneration} and
\ref{prop:root-stage-nonregular}.

Together these results lead to an annular parity cover, a pro-Kummer inverse
system of finite root stacks, a volume line carrying a Jacobi--Kirillov
bracket, a $K$-equivariant theta local system, and a dominating principalized
presentation.  On these objects the two initial GKT Hamiltonians admit an
ordered scattering completion with proper, locally finite support.

\subsection{The geometric result}

\begin{construction}[Flared compactification]\label{con:flare}
Blow up, in the real-oriented sense, the two boundary components of
$\overline\Acal$ and its four focus strata, and attach to the six resulting
ends outward integral-affine cylinders, here called \emph{flares}, of widths
$4,16,1,1,1,1$.  Write $\overline\Acal^{\mathrm{fl}}$ for the result.
\end{construction}

\begin{theorem}[Annular reconstruction]\label{thm:intro-reconstruction}
The quartic toric degeneration and the first two descendent GKT
Hamiltonians determine the following data.

\begin{enumerate}
\item \textup{(The integral-affine surface.)}  The surface $\overline\Acal$ is a
polarized integral-affine annulus
with tangent lattice $\widetilde M$, four primitive focus--focus
singularities, and boundary shear widths $4$ and $16$.  Its polarization is
the minimal integral multiple $4p^*\varphi$ of the pullback of the quartic
polarization.

\item \textup{(The coefficients.)}  The local slab transitions lift after
adjoining roots of degrees $2,4,4,2$.  Closing these roots under transport
gives the inverse system \eqref{eq:intro-root-system}; paths between chambers
act on its finite stages by the lifted transition maps.
For every $N$ and every compact $C\subset\Acal^\circ$, the transported wall
functions over $C$ modulo $I^N$ lie in a single finite root-stack presentation; no finite
root stack contains the complete transported root orbit.

\item \textup{(The wall structure.)}  The inverse system
\eqref{eq:intro-root-system} carries the volume line $\Lcal_\kappa$ and its
Jacobi bracket.  Starting with the two rays in directions $F(2,1)$ and
$F(1,2)$ and the Hamiltonians $\mathsf A$ and $\mathsf B$, the scattering
algorithm gives a unique proper, locally finite wall structure on
$\Acal^\circ$.  Its path-ordered product around every sufficiently small
contractible loop about a joint is the identity.

\item \textup{(Theta multiplication.)}  Given any finite list of theta
sections based in chosen chambers, products among them, associativity
comparisons, and any $N$, there is a finite set of exponents containing the
inputs and every exponent generated in those computations modulo $I^N$.
There is also one finite root-stack stage on which all the corresponding
broken-line expansions are defined and unique.  The resulting products are
commutative and associative.  The core
holonomy $K$ of \eqref{eq:intro-K} acts trivially on $D$ and
hyperbolically on $\widetilde M$.  Consequently
the theta functions form a $K$-equivariant local system; there is no
single-valued basis indexed by a fixed copy of $\widetilde M$.

\item \textup{(Logarithmic compactification.)}  Let
$\overline\Acal^{\mathrm{fl}}$ be as in Construction~\ref{con:flare}.  At every
finite root, coefficient, and deformation stage, transported toric monoid
charts glued along their full boundary overlaps give a separated
fine-saturated logarithmic compactification whose six boundary divisors are
flat over the coefficient ring.  These compactifications commute with the
transition maps of \eqref{eq:intro-root-system}.  Logarithmic regularity fails
at some finite root-stack presentation
(Proposition~\ref{prop:root-stage-nonregular}); the system is nevertheless
dominated by a cofiltered presentation by log-smooth Deligne--Mumford stacks
with regular, log-regular geometric fibers
(Theorem~\ref{thm:regular-presentation}), obtained by functorial logarithmic
resolution, principalization of the pulled-back branch ideals, and fourth
roots of the resulting normal-crossings components.
\end{enumerate}
\end{theorem}

In the orientation of Appendix~\ref{app:calculations}, the two incoming
coefficients in Part~(3) are $-1/2,-1/2$ and the first outgoing coefficient is
$3/4$; these are computed in Section~\ref{sec:jacobi-theta}.

Consistency in Part~(3) has its standard local meaning: the path-ordered
product is the identity around every sufficiently small contractible loop
about a joint.  It does not trivialize the prescribed affine, root, or line
holonomy around a focus or around the annular core.  Part~(5) asserts
separatedness, fine-saturated logarithmic charts, and flatness of the
horizontal boundary at each finite stage.  It asserts neither regularity of
the underlying local rings nor logarithmic regularity, and the inverse limit
is not claimed to be an algebraic stack of finite type.  The final assertion
of Part~(5) concerns the separate dominating presentation of
Theorem~\ref{thm:regular-presentation}; it does not retroactively make the
minimal finite root-stack presentations regular.

The compactification addresses a boundary obstruction absent from the open
locus.  The off-diagonal walls meet the compact boundary transversely,
so the usual convexity argument for tangent boundary walls does not apply.
Nor can the boundary be doubled affinely: its shear is nonzero.  We instead
attach flaring cylinders, order the exiting supports, and transport the
compactified monoid chart across each wall.  Gluing adjacent transported
charts along the full boundary overlap, rather than only after inverting a
boundary parameter, is what gives separatedness and a horizontal flat
divisor.

\subsection{The open-FJRW comparison}

The finite sector group $D$ remembers fractional exponents but not the
ordered pair of quartic twists.  The latter is
\[
 Q_{\GKT}=\widetilde M/4\ZZ^2\simeq(\ZZ/4)^2.
\]
Let $\Acal_{\mathrm{lab}}^\circ\to\Acal^\circ$ be the degree-four cover on
which this quotient local system is constant.  We call it the
\emph{labelled-sector cover}.  A degree-two intermediate cover trivializes
$D$ and has a deck involution induced by interchanging $x$ and $y$.
Coordinate symmetrization below means taking invariants under this
involution after multiplication, followed by the finite-cover pushforward.

The comparison has two logically distinct parts.  First, the $\kappa$-shift
identifies the GKT Lie algebra of critical-graph wall operations with a
labelled subalgebra of the Jacobi algebra of wall Hamiltonians.  A \emph{GKT
endpoint factor} is the ordered product that changes the GKT canonical
boundary condition across one ray.  Dividing the scattering correction on
that ray by this embedded factor defines the \emph{residual factor}.  The
first mixed GKT factor has coefficient $1/4$, while scattering consistency
requires coefficient $3/4$; the residual factor therefore begins with
coefficient $1/2$.  It cannot be omitted without destroying consistency
around the joint.

The Frobenius comparison takes the following second input.

\begin{importedthm}[Primary open mirror theorem; {\cite[Thm.~0.3 and
Cor.~4.17]{grossetal2022open}}]\label{thm:gkt-mirror}
Let $W^{\GKT}_{\mathrm{prim}}$ be the coordinate-symmetric primary open FJRW
potential of $(x^4+y^4,\mu_4^2)$, obtained by setting every descendent
variable to zero.  Then $W^{\GKT}_{\mathrm{prim}}$ is an unfolding of
$x^4+y^4$ with primitive form $dX\wedge dY$ and nine good-basis cycles
$\Xi_{a,b}$, and its nine parameters $t_{a,b}$, $0\leq a,b\leq2$, are
simultaneously Saito-flat coordinates for that unfolding.
\end{importedthm}

What this paper adds to the primary mirror theorem above is the transport,
not the Frobenius structure itself.

Second, a Frobenius comparison requires more than wall coefficients.  It
requires the primitive form, oscillatory periods, a residue pairing, and a
Kodaira--Spencer basis.  These data remove a unipotent ambiguity invisible to
the unary walls.  In Part~(3) below, closure under dependencies means inclusion
of every connected-component datum and every zero- or exchange-boundary datum
required by the GKT assembling relation; Appendix~\ref{app:gkt-completion}
gives the formal definition.  The resulting statement is as follows.

\begin{theorem}[Wall factorization and transported primary structure]
\label{thm:intro-comparison}
On the labelled-sector cover $\Acal_{\mathrm{lab}}^\circ$ and its finite
root-stack stages:

\begin{enumerate}
\item Every finite labelled GKT critical-graph Lie algebra embeds without
changing its coefficients in the Jacobi algebra of wall Hamiltonians, and
every finite GKT endpoint operation maps to the corresponding slope-ordered
product of Jacobi wall automorphisms.  These embeddings commute with
relabelling, restriction, and transport along a framed path.

\item Each total consistency factor has a canonical raywise factorization
into a GKT endpoint factor and a residual unipotent factor.  The residual
factors are compatible under products and changes of framed paths.  For residual
factors $e^\xi,e^\eta,e^\zeta$ of equal deformation degree
the first nontrivial associator has logarithm
\[
 \tfrac1{72}\bigl[[\xi,\zeta],\eta\bigr].
\]

\item Finite GKT endpoint homotopies can be chosen compatibly under
enlargement of finite labelled systems closed under the boundary and
component dependencies of Appendix~\ref{app:gkt-completion}.  Their union is a
componentwise smooth, slope-ordered GKT homotopy whose restriction to every
finite quotient realizes the corresponding endpoint factors.

\item Let $W^{\GKT}_{\mathrm{prim}}$ be the coordinate-symmetric primary GKT
potential and let $P_\Gamma$ be transport along a framed annular path $\Gamma$.  In the
transported Jacobian algebra, the nine classes
\[
 \Phi_{a,b}^{\Gamma}
 =\left[P_\Gamma\!\left(
   \frac{\partial W^{\GKT}_{\mathrm{prim}}}{\partial t_{a,b}}
  \right)\right],
 \qquad 0\leq a,b\leq2,
\]
form a basis.  Transport along another framed path carries this basis to the
corresponding basis there.  Their product, unit, and normalized residue
pairing have structure constants independent of $\Gamma$ in these
transported bases.  Taking invariants under $x\leftrightarrow y$ and pushing
forward along the degree-two intermediate cover gives the complete
narrow-primary GKT/Saito--Givental Frobenius structure in every annular
chart.
\end{enumerate}
\end{theorem}

The last part uses the primary mirror theorem of
\cite{grossetal2022open}.  The contribution here is to place the potential,
primitive form, Brieskorn complex, and Kodaira--Spencer basis in the annular
theta local system and to prove that transport and finite Kummer refinement
preserve the Frobenius tensor.  Section~\ref{sec:primary-frobenius} also checks
the marginal period normalization and distinguishes its coefficients from the
mixed Jacobi wall coefficient $3/4$.

\subsection{Main result: the annular cotangent correspondence}

The positive-descendent comparison is the principal result of the paper and
is independent of the primary transport argument.  In standard terms, its
input is a finite full subdiagram of the stratified quartic $W$-spin moduli
problem, closed under the normalization, detachment, stabilization, and
exchange maps that occur on its boundary.  Its output is a class in relative
oriented PL bordism for a compactified tropical incidence problem.  The next
two definitions make this finite diagram and bordism group precise before the
theorem is stated.

\begin{definition}[Finite labelled window]
\label{def:intro-finite-window}
Each label $i$ has a quartic twist $(a_i,b_i)\in\{0,1,2,3\}^2$ and a
descendent degree $d_i\in\NN$.  An \emph{exact open $W$-spin datum} is the
following finite diagram:
\begin{equation}
 (\Gamma,\mathcal{PM}_\Gamma,\mathbb E_\Gamma(\mathbf d),
  \mathcal L_i,t_{i,w},F_{\Lambda\Gamma},\phi_{\Lambda\Gamma}).
 \label{eq:intro-exact-open-datum}
\end{equation}
Here $\Gamma$ is a stable labelled open quartic $W$-spin graph;
$\mathcal{PM}_\Gamma$ is its compact oriented orbifold with corners;
$\mathbb E_\Gamma(\mathbf d)$ is its descendent Witten bundle;
$\mathcal L_i=T^*_{z_i}\mathcal C_\Gamma$ is the tautological line; and
$t_{i,w}$ is the standard meromorphic-differential section of $\mathcal L_i$
pointing from $z_i$ to a retained marking $w$.  For every boundary graph
$\Lambda$, the datum includes the normalization, detachment, stabilization,
or exchange map
$F_{\Lambda\Gamma}:\mathcal{PM}_\Lambda\to
\partial_\Lambda\mathcal{PM}_\Gamma$ and the corresponding oriented bundle
map $\phi_{\Lambda\Gamma}$.  These maps are the maps of the compact
$W$-spin moduli problem, not numerical open invariants.  They obey the actual
composition identities on multiple boundary strata.  Whenever $i$ is
retained and its component is not contracted, the datum also contains the
canonical splitting
\begin{equation}
 \mathbb E_\Gamma(\mathbf d+\mathbf e_i)
 \simeq \mathbb E_\Gamma(\mathbf d)\oplus\mathcal L_i.
 \label{eq:intro-cotangent-splitting}
\end{equation}
A closed boundary entry analogously specifies its compact closed $W$-spin
orbifold, Euler bundle and section, and its maps on all boundary strata.
Thus an exact open datum is a finite piece of the usual stratified $W$-spin
moduli problem together with the bundles and tautological sections used in
the Euler-chain construction.

A \emph{finite labelled window} $S=(I,P,\mathcal E)$ consists of a finite
labelled set $I$, a finite divisor-closed set $P$ of deformation monomials
containing $1$, and a finite collection $\mathcal E$ of the diagrams
\eqref{eq:intro-exact-open-datum} and their closed boundary entries.  It is
closed under connected components, boundary graphs, and the bundle
splittings \eqref{eq:intro-cotangent-splitting}.  For each support $J$, write
\begin{equation}
 \sum_{i\in J}a_i=r_J+4\ell_J,\qquad
 \sum_{i\in J}b_i=s_J+4m_J,\qquad
 N_J=\ell_J+m_J-|J|+1+\sum_{i\in J}d_i,
 \label{eq:intro-profile-length}
\end{equation}
where $0\leq r_J,s_J\leq3$.  The window also contains the finite sequence with
exponent profiles
\begin{equation}
 (r_J+4p,s_J+4(N_J-p)),\qquad 0\leq p\leq N_J,
 \label{eq:intro-balanced-profiles}
\end{equation}
when $N_J\geq0$, together with the strata joining consecutive profiles.
These are respectively the balanced and critical profiles in the terminology
of \cite{grossetal2022open}.
The window is invariant
under twist-preserving permutations.  Restriction means literal forgetting
of labels, monomials, and their dependent strata.
Equivalently, the window is the finite full subdiagram on which one fixed
coefficient truncation of the open and closed degeneration formulas is
defined.

Put $\mathfrak m_S=(t_i:i\in I)$ and
\begin{equation}
 R_{S,N}=\QQ[t_i:i\in I]\Big/
 \left(\mathfrak m_S^N+\langle m:m\notin P\rangle\right).
 \label{eq:intro-window-ring}
\end{equation}
Divisor closure makes this a finite Artin ring.  Under unlabelling, each
$t_i$ maps to the singleton parameter of the same twist, so the
$\mathfrak m_S$-filtration is the labelled pullback of the singleton
filtration used in the two-generator scattering problem.

A \emph{framed compact annular corridor} is obtained by cutting the affine
annulus along a support-free transverse arc and retaining the compact
polyhedral strip between stops placed before the first and after the last
wall used modulo $\mathfrak m_S^N$.  Its two boundary arcs are the incoming
and outgoing ends.  Stability suppresses every unmarked bivalent wall-free
vertex.  A \emph{finite root stage} is one finite presentation in
\eqref{eq:intro-root-system} containing every root needed by the slab and wall
maps in the corridor and the fourth root of every closed-component smoothing
parameter.
A \emph{target frame} is a terminal chamber and a path from each cell to it;
its transport map is the ordered composition of the wall and root-chart
chain isomorphisms along that path.  Transport around the core gives $K$;
to compare two finite corridors, both are first included in a common finite
enlargement and then transported.
The corridor is therefore a compact cut-open scattering chamber with its two
limiting boundary conditions retained as relative boundary components; a
target frame is only the standard choice needed to compare transported
coefficients in one chart.

The term \emph{broken-line chain} is used at chain level: its cells parametrize
stable decorated tropical forests together with $W$-spin data, rather than
individual Gross--Siebert broken lines.  The finite chain used below is a
cellular chain of stable decorated tropical
trees mapping to this corridor, fibered over the orbifolds
$\mathcal{PM}_\Gamma$ in \eqref{eq:intro-exact-open-datum}.  Forest cuts carry
elements of the transported relative twisted de~Rham complex defined in
Section~\ref{sec:descendent-correspondence}; these elements retain their
values in the chosen finite root extension.  The \emph{contact moment} of a
zero-dimensional boundary chain is its transported relative residue trace.
All endpoint and degeneration facets remain in the cellular boundary unless
they are paired by an explicit mapping cylinder.
\end{definition}

\begin{definition}[Relative PL bordism module]
\label{def:intro-relative-pl-bordism}
Let $\mathfrak X_{S,N}$ be the finite diagram formed from the compact
$W$-spin orbifolds, annular incidence polyhedra, and boundary maps in a finite
labelled window.  Two compact oriented rational PL branched chains in this
diagram are \emph{relatively PL bordant} if they are the two parameter-end
restrictions of a compact oriented rational PL branched chain over $[0,1]$
whose annular-end restrictions are fixed and whose other exterior faces obey
the same labelled boundary maps.  Coefficients are carried in the transported
relative twisted de~Rham complexes.  We write
$\Omega^{\mathrm{PL}}_{S,N}$ for the resulting $R_{S,N}$-module of bordism
classes.  Thus a class in $\Omega^{\mathrm{PL}}_{S,N}$ retains its named
boundary strata; it is not obtained by discarding the degeneration faces.
Section~\ref{sec:descendent-correspondence} gives the cellular realization of
this module.
\end{definition}

\begin{theorem}[Annular cotangent correspondence]
\label{thm:intro-descendent}
Fix a finite labelled window $S=(I,P,\mathcal E)$ as in
Definition~\ref{def:intro-finite-window}, a deformation order $N$, the
finite root stage just defined, and a framed compact annular corridor.  There
is a finite cellular complex $\mathsf C^{\mathrm{dp}}_{S,N}$ over $R_{S,N}$
and a canonically defined relative bordism class
\begin{equation}
 [\BL_{S,N}]_{\mathrm{PL}}\in
 \Omega^{\mathrm{PL}}_{S,N}.
 \label{eq:intro-BL-cobordism-class}
\end{equation}
Its contact moment $\mathcal A_{\mathrm{ann}}(I,\mathbf d)$ satisfies the
following two identities.  For $A\subseteq I$ with $N_A=-1$, let
$Q_A=(I\setminus A)\sqcup\{z_A\}$ be the continuing open datum obtained by
replacing $A$ by its half-node, and let $\mathcal C_A^{\mathrm{ext}}$ be the
Frobenius-normalized Euler number of the complementary closed $W$-spin
factor.  If $j_1,j_2\in I$ are distinct, then
\begin{equation}
\begin{aligned}
 \mathcal A_{\mathrm{ann}}(I,\mathbf d+\mathbf e_{j_1})
 ={}&\sum_{\substack{A\subseteq I,\ j_1\in A,\ |A|\geq2\\N_A=-1}}
 \mathcal C_A^{\mathrm{ext}}
 \mathcal A_{\mathrm{ann}}(Q_A,\mathbf d|_{I\setminus A})
 -\mathcal A_{\mathrm{ann}}(I,\mathbf d),
 \label{eq:intro-qone}
\end{aligned}
\end{equation}
and
\begin{equation}
\begin{aligned}
 \mathcal A_{\mathrm{ann}}(I,\mathbf d+\mathbf e_{j_1})
 ={}&\sum_{\substack{A\subseteq I,\ j_1\in A,\ j_2\notin A,\ |A|\geq2\\
                         N_A=-1}}
 \mathcal C_A^{\mathrm{ext}}
 \mathcal A_{\mathrm{ann}}(Q_A,\mathbf d|_{I\setminus A}).
 \label{eq:intro-qtwo}
\end{aligned}
\end{equation}
Together with
\begin{equation}
 \mathcal A_{\mathrm{ann}}(\varnothing,0)=1,
 \qquad
 \mathcal A_{\mathrm{ann}}(\{i\},d)=(-1)^d,
 \label{eq:intro-descendent-initial-values}
\end{equation}
these identities determine every moment by induction on $|I|$ and identify
the resulting family with the chamber-independent Gross--Kelly--Tessler
moment family.

The class has the following geometric and functorial properties.

\begin{enumerate}
\item It is represented by compact rational PL zero loci over the open and
closed $W$-spin orbifolds in the window.  The Witten-bundle splitting
\eqref{eq:intro-cotangent-splitting} identifies the upper cotangent face with
descendent degree $\mathbf d+\mathbf e_i$; the two annular ends are separate
relative restrictions.  Its remaining exterior boundary
consists of closed-component splittings, wall-crossing homotopies,
cotangent-line contraction faces, the two cyclic orders at a pair collision,
and closed factors whose Euler rank does not match the dimension of their
moduli space.

\item At the two annular endpoints, the coefficients equal the two GKT
boundary systems obtained by retaining, recursively on support, only the
least-slope or only the greatest-slope balanced graph.  Divided-power
pushforward gives the standard automorphism factors.  No arbitrary-chamber
coefficient equality is asserted.

\item The bordism classes commute with window restriction,
deformation-order reduction, finite-root refinement, relabelling, transport
to another target frame, coordinate transposition, and core transport after
including both windows in a common finite stage.  In the inverse limit they
form a $K$-equivariant inverse system of theta-decorated relative PL bordism
classes.  Its coefficientwise contact moment is the complete formal
genus-zero descendent potential.  No map from relative PL bordism to the
theta algebra is asserted.
\end{enumerate}

The geometric moduli spaces, bundles, tautological sections, boundary maps in
\eqref{eq:intro-exact-open-datum}, and the two normalized singleton
Hamiltonians $\mathsf A,\mathsf B$ are inputs.  No higher-support numerical
GKT open invariant, GKT choice of Euler chain, period coefficient, or target
endpoint coefficient enters the construction of
\eqref{eq:intro-BL-cobordism-class}.
\end{theorem}

Section~\ref{sec:descendent-correspondence} gives the construction and proof.
It is finite-dimensional and coefficientwise; it does not produce an
analytic infinite-window virtual fundamental class.

\subsection{Relation with reconstruction and scope}

The construction follows the usual separation between affine and wall data
\cite{grosshackingkeel2015,grossetal2018canonical,grosssiebert2022} but does
not produce a distinguished Carl--Pumperla--Siebert degeneration of the
original quartic pair.  Their
connected anticanonical construction has a disk-shaped intersection complex;
our descendent lattice forces an annular cover with different coarse
topology, and Proposition~\ref{prop:no-CPS-degeneration} shows the two cannot
be reconciled inside that class.  The descendent theorem identifies the
chamber-independent moment family and the coefficients in two specified
extreme normal forms.  It does not identify an unreexpanded pair-of-broken-lines
coefficient or a raw coefficient in an arbitrary GKT chamber.  Those
quantities are not determined by the invariant moments.

The proof has four stages.  First, the discrete Legendre transform and the
descendent exponents determine the enlarged lattice, parity cover, and
pathwise root stacks.  Second, the shifted Hamiltonian bracket places the
initial factors on the rays determined by the cofactor map; filtered scattering and broken-line
re-expansion then give the framed theta algebra, while transported monoid
charts compactify its ends.  Third, the GKT critical algebra is embedded in
the same Jacobi group.  The compatible-homotopy theorem realizes its finite
factors, and the GKT primary mirror theorem supplies the source Frobenius
tensor that is transported through the annular charts.  The fourth stage is
independent of that imported primary tensor.  For each compact open $W$-spin
orbifold in the finite window, we take the fiber product with the finite
incidence space of stable decorated tropical trees in the cut-open annulus.
The equation bundle contains both the open Witten bundle and the annular
incidence equations.  The splitting
$\mathbb E_\Gamma(\mathbf d+\mathbf e_i)=
\mathbb E_\Gamma(\mathbf d)\oplus\mathcal L_i$ makes the added cotangent
power geometric: a PL homotopy in $\mathcal L_i$ joins a generic compatible
section to a signed combination of the standard pointing sections.  The
zeros of those pointing sections are precisely the open--closed boundary
divisors.  Compatible equivariant PL multisections regularize the remaining
equations.  Explicit mapping cylinders fill zero-length tree edges,
collisions of directions, wall homotopies, and the fourth-root normal to an
open--closed divisor.  The two choices of pointing combination give the two
standard one- and two-marking cotangent-line relations.  A relative Fermat
residue trace turns their cellular boundary formulas into scalar recursions,
which prove uniqueness before comparison with the GKT invariants.

Section~\ref{sec:affine-kummer} constructs the affine annulus and its
coefficient system.  Section~\ref{sec:jacobi-theta} proves Jacobi consistency
and theta multiplication.  Section~\ref{sec:compactification} constructs the
logarithmic caps.  Section~\ref{sec:gkt-decoration} gives the labelled GKT
factorization and residual coherence, and Section~\ref{sec:primary-frobenius}
proves the primary comparison.  Section~\ref{sec:descendent-correspondence}
constructs the virtual broken-line chain and proves the descendent
correspondence.  Appendix~\ref{app:calculations} records the
nonstandard affine and logarithmic calculations.  Appendix~\ref{app:gkt-completion}
proves the compatible GKT completion theorem.

The manuscript sources for this paper and its two companions are available at
\url{https://github.com/bwuebben/open-fjrw-scattering}.
